\section{Noise Covariance, Averaging, and Intrinsic Statistical Efficiency}
\label{app:noise-averaging}
\subsection{An exactly solvable stochastic tangent model}
\begin{definition}
In a $g_\star$-orthonormal tangent frame, let $\mathsf H\succ0$ and
consider the quadratic model
\begin{align}
F(x)&=\tfrac12x^\top\mathsf Hx,\\
m_{t+1}&=\beta m_t+(1-\beta)(\mathsf Hx_t+\zeta_t),\\
x_{t+1}&=x_t-\eta m_{t+1},&
\E\zeta_t&=0,&
\E\zeta_t\zeta_t^\top&=\mathsf C .
\label{eq:noise-linear-model}
\end{align}
The noise is IID and independent of the initial state, with finite
second moment. Here $\eta,\beta$ are constant and satisfy
Eq.~\eqref{eq:local-schur}. This model is the linearization of the
stochastic quotient iteration, not an exact global representation
of a nonlinear network loss.
\end{definition}

\begin{theorem}[Discrete Lyapunov covariance]
\label{thm:noise-lyapunov}
With $\alpha=\eta(1-\beta)$ and
\begin{equation}
\mathsf N=\begin{pmatrix}-\alpha I_D\\(1-\beta)I_D\end{pmatrix},
\end{equation}
the centered state covariance converges to the unique positive
semidefinite solution
\begin{align}
\mathsf S&=\mathsf J\mathsf S\mathsf J^\top+
\mathsf N\mathsf C\mathsf N^\top,\\
\mathsf S&=\sum_{k=0}^\infty
\mathsf J^k\mathsf N\mathsf C\mathsf N^\top(\mathsf J^k)^\top,\\
\operatorname{vec}\mathsf S
&=(I-\mathsf J\otimes\mathsf J)^{-1}
\operatorname{vec}(\mathsf N\mathsf C\mathsf N^\top).
\label{eq:noise-lyapunov}
\end{align}
\end{theorem}
\begin{proof}
For $s_t=(x_t,m_t)$,
\begin{align}
s_T&=\mathsf J^Ts_0+\sum_{t=0}^{T-1}
\mathsf J^{T-1-t}\mathsf N\zeta_t,\\
\operatorname{Cov}(s_T)
&=\mathsf J^T\operatorname{Cov}(s_0)(\mathsf J^T)^\top
+\sum_{k=0}^{T-1}
\mathsf J^k\mathsf N\mathsf C\mathsf N^\top(\mathsf J^k)^\top .
\end{align}
The initial term vanishes because $\rho(\mathsf J)<1$.
Any difference of two stationary solutions satisfies
$\Delta=\mathsf J^k\Delta(\mathsf J^k)^\top$ for all $k$, hence
$\Delta=0$. Vectorization gives the final equation since
$\rho(\mathsf J\otimes\mathsf J)=\rho(\mathsf J)^2<1$.
\end{proof}

\begin{theorem}[Exact modal variance and stationary excess risk]
\label{thm:noise-modal-variance}
Diagonalize $\mathsf H$ and put
$\nu_i^2=e_i^\top\mathsf Ce_i$. At stationarity,
\begin{align}
\E x_i^2
&=\frac{\eta(1+\beta)\nu_i^2}
{\lambda_i\{2(1+\beta)-\eta(1-\beta)\lambda_i\}},\\
\E m_i^2
&=\frac{2(1-\beta)\nu_i^2}
{2(1+\beta)-\eta(1-\beta)\lambda_i},\\
\E[x_im_i]
&=-\frac{(1-\beta)\lambda_i}{1+\beta}\E x_i^2,\\
\E F(x)
&=\frac{\eta}{2}\sum_{i=1}^D
\frac{(1+\beta)\nu_i^2}
{2(1+\beta)-\eta(1-\beta)\lambda_i}.
\label{eq:noise-modal-variance}
\end{align}
The noise covariance need not commute with $\mathsf H$.
\end{theorem}
\begin{proof}
Eliminate $m$ in a single mode:
\begin{align}
x_{t+1}&=a_ix_t-\beta x_{t-1}-\alpha\zeta_{t,i},&
a_i&=1+\beta-\alpha\lambda_i .
\end{align}
Write $s_i=\E x_{t,i}^2$ and $c_i=\E[x_{t,i}x_{t-1,i}]$.
Independence of the new noise gives
\begin{align}
c_i&=a_is_i-\beta c_i
\Longrightarrow c_i=\frac{a_i}{1+\beta}s_i,\\
s_i&=(a_i^2+\beta^2)s_i-2a_i\beta c_i+\alpha^2\nu_i^2,\\
s_i\frac{1-\beta}{1+\beta}
\{(1+\beta)^2-a_i^2\}&=\alpha^2\nu_i^2 .
\end{align}
Use
\begin{equation}
(1+\beta)^2-a_i^2
=\alpha\lambda_i\{2(1+\beta)-\alpha\lambda_i\}
\end{equation}
and $\alpha=\eta(1-\beta)$ to obtain the first identity.
Since $m_t=(x_{t-1}-x_t)/\eta$,
\begin{align}
\E m_{t,i}^2&=2(s_i-c_i)/\eta^2,&
\E[x_{t,i}m_{t,i}]&=(c_i-s_i)/\eta .
\end{align}
Finally $F=\frac12\sum_i\lambda_i x_i^2$ uses only diagonal modal
variances, regardless of cross-mode noise correlations.
\end{proof}

\begin{corollary}[The rate--variance distinction]
For fixed $\beta<1$ and $\eta\downarrow0$,
\begin{equation}
\E F(x)
=\frac{\eta}{4}\operatorname{tr}\mathsf C
+\frac{\eta^2(1-\beta)}{8(1+\beta)}
\operatorname{tr}(\mathsf H\mathsf C)+O(\eta^3).
\label{eq:noise-small-step}
\end{equation}
For any mode with $\nu_i^2>0$, its variance diverges as
$\eta(1-\beta)\lambda_i\uparrow2(1+\beta)$.
\end{corollary}
\begin{proof}
Expand each denominator in Eq.~\eqref{eq:noise-modal-variance} by
$(1-u)^{-1}=1+u+O(u^2)$. The strict stability margin makes this
expansion uniform over the finitely many modes.
The divergence follows directly from the exact positive numerator.
\end{proof}

\subsection{Averaging removes the fixed-step modal variance inflation}
\begin{theorem}[Exact boundary identity for the average]
\label{thm:noise-average-boundary}
For $\bar x_T=T^{-1}\sum_{t=0}^{T-1}x_t$ in
Eq.~\eqref{eq:noise-linear-model},
\begin{equation}
\sqrt T\,\bar x_T
=-\mathsf H^{-1}\frac1{\sqrt T}\sum_{t=0}^{T-1}\zeta_t
+\frac{\mathsf H^{-1}}{\sqrt T}
\left\{\frac{x_0-x_T}{\eta}
+\frac{\beta}{1-\beta}(m_T-m_0)\right\}.
\label{eq:noise-average-boundary}
\end{equation}
For square-integrable $s_0$,
\begin{align}
\sqrt T\,\bar x_T
&\ \xrightarrow{\ d\ }\ 
N(0,\mathsf H^{-1}\mathsf C\mathsf H^{-1}),\\
T\,\E[\bar x_T\bar x_T^\top]
&\longrightarrow\mathsf H^{-1}\mathsf C\mathsf H^{-1},\\
T\,\E F(\bar x_T)&\longrightarrow
\tfrac12\operatorname{tr}(\mathsf C\mathsf H^{-1}).
\label{eq:noise-average-covariance}
\end{align}
\end{theorem}
\begin{proof}
Sum both update equations:
\begin{align}
\sum_{t<T}m_{t+1}&=(x_0-x_T)/\eta,\\
(1-\beta)\mathsf H\sum_{t<T}x_t
&=\sum_{t<T}(m_{t+1}-\beta m_t)
-(1-\beta)\sum_{t<T}\zeta_t\\
&=(1-\beta)\sum_{t<T}m_{t+1}
+\beta(m_T-m_0)-(1-\beta)\sum_{t<T}\zeta_t .
\end{align}
This proves the exact identity.
Schur stability and the covariance series imply
$\sup_T\E(\|x_T\|^2+\|m_T\|^2)<\infty$; the boundary term is
therefore $O_{L^2}(T^{-1/2})$.
The IID central limit theorem gives the distributional limit.
The leading noise average has exactly covariance $\mathsf C$;
Cauchy--Schwarz bounds its cross term with the boundary by
$O(T^{-1/2})$. Taking the quadratic-risk trace proves the final claim.
This is an averaging statement, not a claim that the unaveraged
stationary covariance is independent of momentum.
\end{proof}

\subsection{A nonlinear quotient averaging theorem with declared moment stability}
The next result gives the retraction and transport remainder
calculation underlying stochastic averaging
\cite{polyak1992averaging}.

\begin{assumption}[Local stochastic approximation regime]
\label{ass:noise-nonlinear}
In one fixed regular chart around a population minimizer, let
$\mathsf H=\nabla^2(R\circ\psi)(0)\succ0$,
$0\leq\beta<1$ be fixed, and
$\eta_t=c(t+1)^{-\gamma}$ with $1/2<\gamma<1$.
The coordinate state satisfies
\begin{align}
m_{t+1}
&=\beta m_t+(1-\beta)(\mathsf Hx_t+\zeta_{t+1})+r^m_{t+1},\\
x_{t+1}&=x_t-\eta_tm_{t+1}+r^x_{t+1},\\
\E_t\zeta_{t+1}&=0,&
\frac1T\sum_{t<T}\E_t[\zeta_{t+1}\zeta_{t+1}^\top]
&\xrightarrow{\ p\ }\mathsf C,\\
\sup_t\E\|\zeta_{t+1}\|^{2+\delta}&<\infty
\quad\text{for some }\delta>0,\\
\E\|x_t\|^2&\leq C\eta_t,&
\sup_t\E\|m_t\|^2&\leq C,\\
\E\|r^m_{t+1}\|&\leq C\eta_t,&
\E\|r^x_{t+1}\|&\leq C\eta_t^2 .
\label{eq:noise-nonlinear-premises}
\end{align}
Iterates stay in the regular chart, or a localization event with
probability tending to one is used. The displayed last-iterate
moment stability is a premise; it is stronger than an averaged
stationarity estimate.
\end{assumption}

\begin{proposition}[Where the geometric remainders enter]
For the exact quotient iteration, the last two remainder bounds in
Eq.~\eqref{eq:noise-nonlinear-premises} follow locally from the
displayed moment bounds if the gradient has a quadratic Taylor
remainder, retraction and transport have bounded local derivatives,
and the pre-transport direction has uniformly bounded second moment.
\end{proposition}
\begin{proof}
For constants uniform on the chart,
\begin{align}
g(x)&=\mathsf Hx+\mathcal R_g(x),&
\|\mathcal R_g(x)\|&\leq C_g\|x\|^2,\\
[\operatorname{Retr}_X(-\eta v)]_\psi
&=x-\eta v+\mathcal R_R(x,v,\eta),&
\|\mathcal R_R\|&\leq C_R\eta^2\|v\|^2,\\
[\mathcal T_{X,X^+}]_{E(X),E(X^+)}
&=I+\mathcal R_T,&
\|\mathcal R_T\|_{\rm op}&\leq C_T\eta\|v\| .
\end{align}
With $\widetilde m=\beta m+(1-\beta)(g(x)+\zeta)$ and
$m^+=\mathcal T\widetilde m$,
\begin{align}
r^m&=(1-\beta)\mathcal R_g(x)+
\mathcal R_T\widetilde m,\\
r^x&=\mathcal R_R+\eta(m^+-\widetilde m).
\end{align}
Cauchy--Schwarz, $\E\|x\|^2=O(\eta)$, and the second-moment
bounds give $\E\|r^m\|=O(\eta)$ and
$\E\|r^x\|=O(\eta^2)$. All frame and transport corrections are
retained in these remainders.
\end{proof}

\begin{theorem}[Asymptotic linearity of averaged quotient iterates]
\label{thm:noise-nonlinear-average}
Under Assumption~\ref{ass:noise-nonlinear},
\begin{align}
\sqrt T\,\bar x_T
&=-\mathsf H^{-1}\frac1{\sqrt T}
\sum_{t<T}\zeta_{t+1}+o_p(1),\\
\sqrt T\,\bar x_T
&\xrightarrow{\ d\ }
N(0,\mathsf H^{-1}\mathsf C\mathsf H^{-1}).
\label{eq:noise-nonlinear-average}
\end{align}
If, in addition, $\sup_T\E\|\sqrt T\bar x_T\|^{2+\delta}<\infty$
and the population Taylor remainder is
$O(\|x\|^{2+\delta})$ for some $0<\delta\leq1$, then
\begin{equation}
T\,\E[R(\psi(\bar x_T))-R(X_\star)]
\longrightarrow\tfrac12\operatorname{tr}(\mathsf C\mathsf H^{-1}).
\label{eq:noise-nonlinear-risk}
\end{equation}
\end{theorem}
\begin{proof}
Eliminate $m_{t+1}$ in the two recursions and sum:
\begin{align}
\mathsf H\sum_{t<T}x_t
&=-\sum_{t<T}\zeta_{t+1}
+\sum_{t<T}\frac{x_t-x_{t+1}}{\eta_t}
+\frac{\beta}{1-\beta}(m_T-m_0)\notag\\
&\quad+\sum_{t<T}\frac{r^x_{t+1}}{\eta_t}
-\frac1{1-\beta}\sum_{t<T}r^m_{t+1}.
\label{eq:noise-abel-identity}
\end{align}
Abel summation gives the exact boundary expansion
\begin{align}
\sum_{t<T}\frac{x_t-x_{t+1}}{\eta_t}
&=\frac{x_0}{\eta_0}-\frac{x_T}{\eta_{T-1}}
+\sum_{t=1}^{T-1}x_t
\left(\frac1{\eta_t}-\frac1{\eta_{t-1}}\right),\\
\E\left\|\sum_{t<T}\frac{x_t-x_{t+1}}{\eta_t}\right\|
&\leq C(1+T^{\gamma/2}),\\
\E\sum_{t<T}
\left(\|r^x_{t+1}\|/\eta_t+\|r^m_{t+1}\|\right)
&\leq C\sum_{t<T}\eta_t=O(T^{1-\gamma}).
\end{align}
After division by $\sqrt T$, these errors are
$O(T^{(\gamma-1)/2}+T^{1/2-\gamma})=o(1)$.
The moment bound also makes the momentum endpoint term $o_{L^1}(1)$.
Uniform $(2+\delta)$ moments give the conditional Lindeberg condition:
\begin{equation}
\frac1T\sum_{t<T}\E_t\left[
\|\zeta_{t+1}\|^2\mathbf1_{\{\|\zeta_{t+1}\|>\epsilon\sqrt T\}}
\right]\xrightarrow{\ p\ }0 .
\end{equation}
The martingale central limit theorem and the conditional covariance
premise prove the Gaussian limit. For the final claim, uniform
$(2+\delta)$ moments allow passage of quadratic expectations, while
\begin{equation}
T\,\E\|\bar x_T\|^{2+\delta}
=T^{-\delta/2}\E\|\sqrt T\bar x_T\|^{2+\delta}\to0 .
\end{equation}
Expand the population risk to second order.
\end{proof}

\begin{corollary}[An implementable rank-preserving average]
Let $\bar Y_T=T^{-1}\sum_{t<T}X_t$ and let
$\Pi_r(\bar Y_T)$ be its best rank-$r$ Frobenius approximation.
Under the preceding assumptions,
\begin{equation}
\sqrt T\{\Pi_r(\bar Y_T)-\psi(\bar x_T)\}\xrightarrow{\ p\ }0 .
\label{eq:noise-rank-average}
\end{equation}
Thus its first-order product-space Gaussian limit is the same.
\end{corollary}
\begin{proof}
Near a regular $X_\star$, the spectral gap
$\sigma_r(X_\star)>0=\sigma_{r+1}(X_\star)$ gives a smooth local
nearest-rank projection, with
\begin{equation}
\Pi_r(X_\star+E)=X_\star+P_{X_\star}E+O(\|E\|_F^2).
\end{equation}
No distinctness of the positive singular values is required.
For $J=D\psi(0)$,
\begin{align}
X_t&=X_\star+Jx_t+O(\|x_t\|^2),\\
\bar Y_T&=X_\star+J\bar x_T+
O_p\!\left(T^{-1}\sum_{t<T}\eta_t\right)
=X_\star+J\bar x_T+O_p(T^{-\gamma}).
\end{align}
Since $P_{X_\star}J=J$, the difference in the claim is
$O_p(T^{-\gamma}+\|\bar x_T\|^2)$, which is $o_p(T^{-1/2})$.
\end{proof}

\begin{proposition}[Coordinate and gauge invariance of statistical efficiency]
For $y=Tx+O(\|x\|^2)$ with invertible $T$,
\begin{align}
\mathsf H_y&=T^{-\top}\mathsf H_xT^{-1},&
\mathsf C_y&=T^{-\top}\mathsf C_xT^{-1},\\
\mathsf H_y^{-1}\mathsf C_y\mathsf H_y^{-1}
&=T(\mathsf H_x^{-1}\mathsf C_x\mathsf H_x^{-1})T^\top,\\
\operatorname{tr}(\mathsf C_y\mathsf H_y^{-1})
&=\operatorname{tr}(\mathsf C_x\mathsf H_x^{-1}).
\end{align}
\end{proposition}
\begin{proof}
At a stationary population point, Hessians transform by congruence;
per-example gradients transform as covectors. Apply the chain rule
and the cyclic trace identity. Gauge representatives induce the
same quotient estimator and therefore the same scalar efficiency
constant.
\end{proof}
